\usepackage{verbatim}
\usepackage{scicite}
\usepackage{todonotes}

\newcommand{\Armando}[1]{\todo[author=Armando]{#1}}
\newcommand{\Kevin}[1]{\todo[author=Kevin]{#1}}

\usepackage{times}

\usepackage{multibib}

%% \topmargin 0.0cm
%% \oddsidemargin 0.2cm
%% \textwidth 16cm 
%% \textheight 21cm
%% \footskip 1.0cm


\usepackage{dblfloatfix}    % To enable figures at the bottom of page
% Recommended, but optional, packages for figures and better typesetting:
\usepackage{microtype}
\usepackage{graphicx}
%\usepackage{subfigure}
\usepackage{booktabs} \usepackage{stmaryrd}% for professional tables

\usepackage{verbatim}
\usepackage{longtable}
\usepackage[section]{placeins}
% hyperref makes hyperlinks in the resulting PDF.
% If your build breaks (sometimes temporarily if a hyperlink spans a page)
% please comment out the following usepackage line and replace
% \usepackage{icml2018} with \usepackage[nohyperref]{icml2018} above.
%\usepackage{hyperref}

% Attempt to make hyperref and algorithmic work together better:
\newcommand{\theHalgorithm}{\arabic{algorithm}}
\newcommand{\dcaption}[1]{\caption{#1}}
\newcommand{\system}{DreamCoder~}
\newcommand{\systemEnding}{DreamCoder}
\newcommand{\lowerBound}{\mathscr{L}}
\newcommand{\denotation}[1]{{\llbracket #1 \rrbracket}}
\newcommand{\code}[1]{{\footnotesize\texttt{#1}}}
\newcommand{\scode}[1]{{\tiny\texttt{#1}}}
\newcommand{\mcode}[1]{{\scriptsize\texttt{#1}}}
\newcommand{\codechar}[1]{{\footnotesize{\texttt{"#1"}}}}
\newcommand{\remove}[1]{{\color{red}{#1}}}
% Use the following line for the initial blind version submitted for review:
%\usepackage[nohyperref]{icml2018}

\newcommand{\program}{\rho}
\newcommand{\prior}{L}
\newcommand{\library}{D}

\usepackage{xcolor}
\definecolor{pop1}{HTML}{1F78b4}
\definecolor{pop2}{HTML}{164C13}
\definecolor{pop3}{HTML}{d95F02}
\definecolor{orange}{HTML}{d95F02}
\definecolor{teal}{HTML}{1b9e77}
\newcommand{\pop}[1]{\textcolor{pop1}{#1}}
\newcommand{\popp}[1]{\textcolor{pop2}{#1}}
\newcommand{\tree}[1]{\textcolor{pop3}{#1}}
\newcommand{\orange}[1]{\textcolor{orange}{#1}}
\newcommand{\teal}[1]{\textcolor{teal}{#1}}

\newcommand{\greenCode}[1]{{\footnotesize\popp{\code{#1}}}}
\newcommand{\blueCode}[1]{{\footnotesize\pop{\code{#1}}}}

%\usepackage{hyperref}       % hyperlinks
\usepackage{url}            % simple URL typesetting
\usepackage{booktabs}       % professional-quality tables
\usepackage{amsfonts}       % blackboard math symbols
\usepackage{nicefrac}       % compact symbols for 1/2, etc.
\usepackage{microtype}      % microtypography
\usepackage{mathrsfs}
\usepackage{listings}
\usepackage{amsthm}
% use Times
\usepackage{times}
% For figures
\usepackage{graphicx} % more modern
%\usepackage{epsfig} % less modern
%\usepackage{epstopdf}
\usepackage{subfig} 
\usepackage{fancyvrb}


\usepackage{caption}
%\usepackage{subcaption}

\fvset{fontsize=\footnotesize}

\usepackage{amssymb}
\usepackage{listings}
\usepackage{wrapfig}
\usepackage{tabularx}


%\usepackage{verbatim}
 \usepackage{booktabs}
 % For algorithms
\usepackage{algorithm}
%\usepackage{algorithmic}
\usepackage{algpseudocode}% http://ctan.org/pkg/algorithmicx
\usepackage{tikz}
\usepackage{circuitikz}
\usetikzlibrary{fit,bayesnet,calc,tikzmark}
\usetikzlibrary{arrows.meta}
\usetikzlibrary{positioning}
\usetikzlibrary{decorations.text}

\usepackage{dsfont}
\usepackage{amsmath}

\DeclareMathOperator*{\argmin}{arg\,min} % thin space, limits underneath in displays
\DeclareMathOperator*{\argmax}{arg\,max} % thin space, limits underneath in displays
 


% Packages hyperref and algorithmic misbehave sometimes.  We can fix
% this with the following command.

%\newcommand{\Expect}{\mathds{E}} %{{\rm I\kern-.3em E}}
\newcommand{\Expect}{\text{E}} %{{\rm I\kern-.3em E}}
\newcommand{\indicator}{\mathds{1}} %{{\rm I\kern-.3em E}}
%\newcommand{\expect}{\mathds{E}} %{{\rm I\kern-.3em E}}
\newcommand{\expect}{\text{E}} %{{\rm I\kern-.3em E}}
%\newcommand{\probability}{\mathds{P}} %{{\rm I\kern-.3em P}}
\newcommand{\probability}{\text{P}} %{{\rm I\kern-.3em P}}
\newcommand{\shift}[1]{\uparrow^{#1}}
\newcommand{\downshift}[1]{\downarrow^{#1}}
\newcommand{\substitute}[2]{[\$ #1 \mapsto #2]}
\newcommand{\reduce}{\longrightarrow}
\newcommand{\manyReduce}{\longrightarrow^*}

\newtheorem{definition}{Definition}
\newtheorem{theorem}{Theorem}
\newtheorem{lemma}{Lemma}
